\documentclass[12pt]{article}
\usepackage{amsmath,amssymb,amsfonts}
\usepackage[margin=1in]{geometry}

\begin{document}

\section*{Chapter 3, Problem 11}

\textbf{Problem.} Prove that any antisymmetric matrix of odd order is singular. The ``odd order'' qualification is essential; the Pauli spin matrix
\[
\sigma_y = \begin{pmatrix} 0 & -i \\ i & 0 \end{pmatrix}
\]
is a second-order antisymmetric matrix, but since $\sigma_y^2 = I$, $\sigma_y$ does have an inverse, namely $\sigma_y^{-1} = \sigma_y$.

\bigskip
\textbf{Solution.}

Let $A$ be an $n \times n$ antisymmetric (skew-symmetric) matrix with $n$ odd. Then $A^T = -A$.

Recall that for any $n \times n$ matrix, $\det(A^T) = \det(A)$ and $\det(cA) = c^n \det(A)$.

From $A^T = -A$:
\[
\det(A^T) = \det(-A) = (-1)^n \det(A).
\]
Since $\det(A^T) = \det(A)$:
\[
\det(A) = (-1)^n \det(A).
\]

When $n$ is odd, $(-1)^n = -1$, so:
\[
\det(A) = -\det(A) \quad \Longrightarrow \quad 2\det(A) = 0 \quad \Longrightarrow \quad \det(A) = 0.
\]

Since $\det(A) = 0$, the matrix $A$ is singular.

\medskip
\textbf{Why the qualification ``odd order'' is essential:} When $n$ is even, $(-1)^n = 1$, so the equation $\det(A) = \det(A)$ is automatically satisfied and imposes no constraint on $\det(A)$. The example $\sigma_y$ is a $2 \times 2$ antisymmetric matrix with $\det(\sigma_y) = (0)(0) - (-i)(i) = -1 \neq 0$, confirming it is nonsingular. \hfill $\blacksquare$

\end{document}
